\section{Nhận xét và kết luận}

\subsection{Nhận xét tổng hợp}

Đồ án tập trung vào tập \texttt{car\_detail.csv} gồm 30.652 tin đăng xe và 21 cột, phản ánh thị trường xe mới và xe đã dùng tại Việt Nam trong phạm vi năm sản xuất 1990 - 2023. Các thống kê mô tả cho thấy xe đã dùng chiếm tỷ trọng lớn (80,2\%), xe lắp ráp trong nước nhiều hơn xe nhập khẩu (58,3\% so với 41,7\%), hai dòng xe nổi bật nhất là SUV và Sedan, còn Toyota là hãng có số lượng bản ghi cao nhất.

Về kỹ thuật xe, dữ liệu nghiêng mạnh về xe dùng xăng (80,7\%) và hộp số tự động (82,1\%). Số chỗ phổ biến nhất là 5 chỗ, phù hợp với các dòng xe gia đình và đô thị. Dữ liệu cũng tập trung nhiều ở các năm 2018 - 2023, đặc biệt năm 2023 có 4.006 bản ghi, cao nhất trong nhóm năm gần được nêu trong mô tả.

\subsection{Trả lời câu hỏi trung tâm}

Câu hỏi trung tâm của đồ án được xác định là: \textit{``Dữ liệu tin đăng ô tô cho thấy thị trường ô tô Việt Nam được cấu trúc như thế nào về nguồn cung, mức giá và quy luật khấu hao; đồng thời phản ánh ra sao sự chuyển dịch sang xe điện, thị hiếu của người mua và những điểm bất thường cần lưu ý?''}

Dựa trên kết quả của 5 tab phân tích, thị trường trong tập dữ liệu có mức độ tập trung tương đối cao vào một số thương hiệu phổ biến, đặc biệt là Toyota, Hyundai, Ford, Mercedes-Benz và Kia. SUV và Sedan là hai kiểu dáng chủ đạo; xe đã qua sử dụng chiếm phần lớn nguồn cung, trong khi xe lắp ráp trong nước nhỉnh hơn xe nhập khẩu. Mức giá trung vị 638 triệu VND cho thấy phân khúc phổ thông và trung cấp là vùng thị trường trung tâm, mặc dù vẫn tồn tại một nhóm nhỏ xe siêu sang tạo ra khoảng cách rất lớn giữa hai đầu phân phối giá.

Giá trị phương tiện chịu tác động rõ rệt của tuổi đời và quãng đường sử dụng. Giai đoạn khoảng năm năm đầu và ngưỡng 100.000 km là những mốc đáng chú ý, sau đó tốc độ giảm giá có xu hướng chậm lại. Vì vậy, việc định giá xe không thể chỉ dựa trên thương hiệu mà cần kết hợp năm sản xuất, phiên bản, tình trạng và cường độ khai thác thực tế.

Ở chiều chuyển dịch công nghệ, xe xăng vẫn chiếm ưu thế tuyệt đối trong toàn bộ dữ liệu, nhưng xe điện và hybrid tăng nhanh trong nhóm xe đời mới. Xe điện có mức giá cạnh tranh và tỷ lệ xe mới cao, cho thấy quá trình chuyển đổi sang phương tiện năng lượng mới đã bắt đầu thể hiện trên nguồn cung. Về thị hiếu, người mua ưu tiên màu trung tính, hộp số tự động và cấu hình 5 chỗ; quãng đường trung vị của xe cũ khi được rao bán có xu hướng giảm, gợi ý chu kỳ sở hữu và nâng đời xe đang rút ngắn.

Bên cạnh các quy luật chung, dữ liệu còn chứa những điểm cực đoan về giá và đột biến nguồn cung, tiêu biểu là biên độ giá hàng chục tỷ đồng trong cùng một thương hiệu và lượng xe Ford sản xuất năm 2023 tăng bất thường. Các trường hợp này cần được xem là tín hiệu phân tích cần xác minh, thay vì tự động loại bỏ như nhiễu. Như vậy, 5 tab trên dashboard đã tạo thành một mạch phân tích thống nhất: đi từ cấu trúc thị trường, cơ chế định giá, chuyển dịch năng lượng và hành vi người mua đến việc nhận diện các ngoại lệ có thể ảnh hưởng đến quyết định kinh doanh và mô hình dữ liệu.

\subsection{Kết luận và hạn chế}

Báo cáo đã xây dựng được một quy trình phân tích tương đối hoàn chỉnh, từ khảo sát và chuẩn hóa dữ liệu đến tổ chức mô hình và trực quan hóa trên Power BI. Năm tab phân tích bổ sung cho nhau, giúp người xem không chỉ nhận biết nhóm xe và thương hiệu phổ biến mà còn đánh giá mặt bằng giá, quy luật khấu hao, tốc độ thâm nhập của xe điện, thị hiếu cấu hình và các điểm bất thường trong dữ liệu. Kết quả có thể hỗ trợ người mua lựa chọn thời điểm và phân khúc phù hợp, đồng thời cung cấp thông tin tham khảo cho đại lý trong quản lý danh mục và kiểm soát rủi ro định giá.

Tuy nhiên, báo cáo vẫn tồn tại một số hạn chế sau:
\begin{itemize}
    \item \textbf{Phạm vi đại diện của dữ liệu:} Tập dữ liệu được thu thập từ các tin đăng trên một nền tảng trực tuyến nên phản ánh cơ cấu nguồn cung được niêm yết trên nền tảng đó, không đại diện đầy đủ cho toàn bộ thị trường ô tô Việt Nam hoặc các giao dịch ngoại tuyến.
    \item \textbf{Giá chào bán không phải giá giao dịch:} Trường giá thể hiện mức người bán niêm yết, chưa phản ánh giá thương lượng cuối cùng, thời gian bán xe, chi phí sang tên hoặc các ưu đãi đi kèm. Vì vậy, các kết luận về định giá nên được hiểu là phân tích giá chào bán.
    \item \textbf{Các điểm bất thường chưa được xác minh bên ngoài:} Những trường hợp như đột biến nguồn cung xe Ford sản xuất năm 2023 hoặc các mức giá cực đoan mới được phát hiện từ dữ liệu nội bộ. Cần đối chiếu với thông tin đại lý, báo cáo ngành hoặc dữ liệu giao dịch thực tế trước khi diễn giải thành biến động của thị trường.
    \item \textbf{Tính thời điểm của dashboard:} Dashboard phản ánh một ảnh chụp dữ liệu tại thời điểm thu thập và chưa được kết nối với nguồn cập nhật tự động. Các xu hướng và tỷ trọng có thể thay đổi khi bổ sung tin đăng mới hoặc mở rộng sang những nền tảng khác.
\end{itemize}

Trong các nghiên cứu tiếp theo, có thể mở rộng dữ liệu theo thời gian và khu vực, bổ sung giá giao dịch thực tế, thông tin phiên bản và lịch sử xe, đồng thời xây dựng các mô hình dự đoán giá và phát hiện bất thường có bước kiểm định độc lập. Những cải tiến này sẽ giúp dashboard chuyển từ công cụ mô tả sang hệ thống hỗ trợ quyết định có độ tin cậy cao hơn.



\clearpage